\section{Results and precise comparison}
\label{sec:results}
\newcommand{\PrincipalVariables}{215}
\newcommand{\PrincipalLast}{214}
\newcommand{\PrincipalRows}{154}
\newcommand{\PrincipalLastRow}{153}
\newcommand{\PrincipalTerms}{7081}
\newcommand{\PrincipalHeight}{8880}

The principal ordinary polynomial is
\begin{equation}\label{eq:principal}
 \boxed{P_{\PrincipalVariables}(x_0,\ldots,x_{\PrincipalLast})
       =\sum_{j=0}^{\PrincipalLastRow}R_j(x)^2=0,
       \qquad x\in\Z^{\PrincipalVariables}.}
\end{equation}
Every residual is printed in Appendix~\ref{sec:residuals}. Its finite
literal array specifies every coefficient and independent coordinate;
no variable exponent, oracle, recurrence, sequence or hidden witness
is an algebraic primitive. The final collected polynomial has exact
degree four, \PrincipalTerms\ monomials and coefficient height
\PrincipalHeight.

\begin{theorem}\label{thm:main}
For the principal polynomial and each selected companion on its stated
domain $D$, the ordinary witness arguments composed with the identified
inherited syntactic calibration prove
\[
 \mathrm{ZFC}\vdash
 [\neg\exists x\in D^n\;P(x)=0]\ \longleftrightarrow\
 \Con(\mathrm{ZFC}).
\]
No consistency assumption is used to prove this equivalence.
\end{theorem}
This does not assert that either side is true. In particular it neither
proves consistency nor decides whether an equation has a zero.

\subsection{Direct quartics over unrestricted integers}
\begin{center}\small
\begin{tabular}{@{}lrrrr@{}}\toprule
Endpoint & Variables & Rows & Terms & Height\\\midrule
Round21 principal &217&154&7396&8880\\
Incoming CRT-anchor principal &215&154&5660&91679\\
Round22 principal &\PrincipalVariables&\PrincipalRows&\PrincipalTerms&\PrincipalHeight\\
Round22 balanced &215&154&7315&4096\\
Round22 lower height &215&154&16077&2304\\
Round22 lower dimension &214&153&7078&1052028\\\bottomrule
\end{tabular}
\end{center}
Every row represents a distinct fixed quartic. The incoming anchor point
is sparser than the selected moderate-height presentation; it is not
claimed to be dominated. The 214-variable endpoint pays a substantially
larger coefficient height, and the lower-height endpoint pays more
support. The balanced equation uses 234 more monomials than the principal
and reduces height from 8880 to 4096; the exact catalog preserves that
tradeoff rather than hiding it in a single simplicity score.

The incoming source contributes the CRT-anchor and label-separation ideas.
The new contributions are their actual composition with the transported
quadratic simplifications, the weighted and quadratic growth budgets,
and independently checked final coefficient arrays. Savings from different
sources are not added without implementing their joint residual system.

\subsection{A separate lower dimension tradeoff}
The reciprocal companion has 212 unrestricted integer variables, 151 residuals,
55820 monomials and coefficient height 31104, but exact degree 10 rather
than four. Its proof in Section~\ref{sec:reciprocal22} uses a cofinal
$q\equiv1\pmod3$ restriction, followed by reciprocal pooling of the
$q$ and $q+1$ certificates. The original even $H$ is unchanged. Its
higher degree and much larger height are explicit costs, not omitted
coordinates in the quartic comparison. Against source 168's separate
degree-ten point in 215 integers with 152 rows, 99914 monomials and
height 113304, it improves all four displayed finite size measures;
its arithmetic front is nevertheless different, not a relabeling of
that source's zero set.

\subsection{Height one and positive scalar equations}
The selected height-one catalog ranges from $493/24828$ to $701/7141$
variables/monomials. Its nine points include the retained $557/11991$
tradeoff. Each is a quartic with constant one and all collected nonzero
coefficients in $\{-1,1\}$, on unrestricted integer coordinates.

The scalar domain is instead exactly
\[
 (J,F,\alpha,i,T,w,s,n,\eta,\zeta)\in\Z_{>0}^{10}.
\]
The three exact degree/live-operation pairs are $214/87$, $232/85$ and
$240/82$, unchanged from Round21. The new complete four-port tuple occupies
97046 minimal unsigned bytes instead of 295388, a 67.146 percent reduction.
This is a raw numerical payload comparison, not compressed ZIP size,
polynomial support or the height of the fully expanded scalar polynomial.
Those latter two scalar measures are not claimed. The separate
11-unrestricted-integer six-port family is unchanged and is not a new
selected artifact in this round.
